% 1. 前置知识：编辑器、编译器、IDE
\section{前置知识：编辑器、编译器、IDE}

\begin{frame}{编辑器、编译器、IDE}
    写代码之前，先分清三样东西。

    \textbf{编辑器}：任何可以编辑文本的软件，例如记事本、Word、WPS文字、手机备忘录。

    但我们真正需要的是\textbf{代码编辑器}，尤其是\textbf{现代代码编辑器}：\code{Notepad++}、\code{Visual Studio Code}、\code{Vim}、\code{Sublime Text}……

    \textbf{编译器}：用来将高级语言的代码翻译成计算机真正能执行的代码。流行的C/C++编译器有\code{GCC}（GNU）、\code{Clang}（LLVM）、\code{MSVC}（Microsoft）、\code{ICC}（Intel）等。

    \textbf{IDE}（Integrated Development Environment，集成开发环境）：代码编辑器$+$编译器$+$调试器$+$项目管理工具$+$……
\end{frame}

\begin{frame}{我们到底需要什么？}
    \begin{itemize}
        \item 好的IDE有各种工具，例如linter、profiler、版本管理等。
        \item \textbf{宇宙最强IDE}：Microsoft Visual Studio。
        \item 常见的IDE：VS、CLion、PyCharm……
        \item 我们实际需要的是一个IDE——但不必是它的“打包形态”。
    \end{itemize}

    两条路线都行：
    \begin{itemize}
        \item 直接用CLion或VS这样的狠货；
        \item 用编辑器$+$编译器$+$插件自己搭一套，这也是本课程的推荐路线。
    \end{itemize}
\end{frame}
